% Limitations --- required by ACL/EMNLP

\paragraph{Architecture and scale.}
Our primary results are on DeepSeek-R1-Distill-Qwen-7B~\citep{deepseek2025r1} with GQA ($G{=}7$).
Cross-model validation on Qwen2.5-Math-7B covers SFT recovery only ($N{=}200$); both models share the Qwen2.5 base architecture, so our cross-model evidence does not extend to different architecture families (e.g., LLaMA, Mistral).
Whether the spectrum generalizes across architecture families and scales (1B--70B) remains open.

\paragraph{Diagnostic limitations.}
Reward-accuracy decoupling at 20--25\% sparsity means that recovery capacity, operationalized as peak reward, does not predict downstream accuracy improvements within our training horizon.
Recovery capacity captures the model's ability to sustain RL training without collapse, but practitioners seeking accuracy gains may require longer training, different reward formulations, or lower sparsity levels.
The per-layer capacity floor (40\% retention) was chosen heuristically; an ablation over $\tau \in \{0.3, 0.4, 0.5, 0.6\}$ (Appendix~\ref{sec:appendix_floor}) suggests moderate sensitivity, but a full search over continuous $\tau$ values was not performed.

\paragraph{Inference efficiency.}
At 20\% sparsity, pruning reduces memory by ${\sim}$28\% and improves time-to-first-token, but inference throughput decreases by 19\% (72.0$\to$58.3 tok/s), likely because channel-group pruning creates irregular sparse patterns that current hardware cannot accelerate.

\paragraph{Replication and scope.}
CKA-only, Standard, and Random are validated over 3 GRPO seeds; \rasp{} at $\alpha{=}1.0$ uses a single seed ($N{=}1$).
Random uses a 512-token generation limit (vs.\ 1{,}024 for other methods), preventing strict reward comparisons; its ranking should be interpreted as a conservative bound.
We evaluate on three benchmarks (GSM8K, MATH-500, MBPP) covering math reasoning and code generation; logical, commonsense, and multi-modal reasoning are not addressed.
